\section{ফর্ম: সহজ ডেটা এন্ট্রি}
\begin{definitionbox}{ফর্ম (Form)}
ব্যবহারকারীর সুবিধামতো বিন্যস্ত পর্দার মাধ্যমে ডেটা দেখা, যোগ বা সম্পাদনার জন্য ব্যবহৃত ডেটাবেজ অবজেক্টকে \textbf{ফর্ম (Form)} বলে।
\end{definitionbox}
কলেজের ভর্তি ফর্মে নাম, জন্মতারিখ ও বিভাগ আলাদা ঘরে লেখা যায়। পেছনে তথ্য টেবিলে সংরক্ষিত হয়; ফর্ম সেই তথ্য ব্যবহারের ইন্টারফেস। তাই ফর্ম তৈরি করলেই একই তথ্যের নতুন স্বতন্ত্র টেবিল তৈরি হয় না।

\subsection{ফর্মের উৎস ও কন্ট্রোল}
ফর্মের Record Source সাধারণত টেবিল বা কুয়েরি। একটি ফিল্ডের সঙ্গে যুক্ত \textbf{Bound Control} তার মান দেখায় বা সম্পাদনা করে। কোনো ফিল্ডের সঙ্গে যুক্ত নয় এমন \textbf{Unbound Control} নির্দেশনা বা ইনপুট গ্রহণের সহায়ক অংশ হতে পারে। হিসাবের expression ব্যবহার করলে calculated control পাওয়া যায়। যেমন Quantity ও UnitPrice থেকে Total দেখানো যায়।

\textbf{Label} পরিচয় বা নির্দেশনা দেখায়; ``শিক্ষার্থীর নাম'' একটি label। \textbf{Text Box} নাম, নম্বর বা অন্য মান প্রদর্শন ও এন্ট্রিতে ব্যবহৃত হয়। label-এর লেখা বদলালে ফিল্ডের সংরক্ষিত নাম নিজে থেকে বদলায় না। \textbf{Combo Box} তালিকা থেকে মান বাছাই করতে দেয়; যেমন বিভাগ নির্বাচন। এতে বানানের ভিন্নতা কমানো যায়, তবে সঠিক তালিকা ও গ্রহণের নিয়ম নির্ধারণ করতে হয়।

\textbf{Check Box} Yes/No ধরনের অবস্থা প্রকাশ করে; যেমন ফি পরিশোধিত কি না। \textbf{Command Button} চাপলে নির্দিষ্ট কাজ হয়, যেমন নতুন রেকর্ডে যাওয়া বা রিপোর্ট খোলা। বোতামের নাম দিলেই কাজ সক্রিয় হয় না; সংশ্লিষ্ট action বা event নির্ধারণ থাকতে হয়।
\begin{flushleft}\begin{minipage}{\linewidth}\raggedright\small
\begin{tabularx}{0.88\linewidth}{|p{34mm}|X|}\hline
\chSixHeader \multicolumn{2}{|l|}{\textbf{Student Entry Form}}\\\hline
শিক্ষার্থীর নাম & Rupa\\\hline
বিভাগ & Science \quad (তালিকা থেকে নির্বাচন)\\\hline
জন্মতারিখ & 15/05/2008\\\hline
বর্তমানে অধ্যয়নরত & হ্যাঁ \quad (Check Box)\\\hline
\multicolumn{2}{|l|}{নতুন রেকর্ড \qquad পরবর্তী রেকর্ড}\\\hline
\end{tabularx}
\captionof{figure}{ফর্মের কন্ট্রোল বোঝানোর ধারণাগত বিন্যাস}
\end{minipage}\end{flushleft}

\subsection{ফর্ম তৈরির পদ্ধতি ও ভিউ}
Access-এ Form নির্দেশ দিয়ে নির্বাচিত উৎস থেকে দ্রুত ফর্ম তৈরি করা যায়। \textbf{Form Wizard} ধাপে ধাপে উৎস, ফিল্ড ও বিন্যাস নির্বাচন করায়। \textbf{Form Design} দিয়ে প্রয়োজনমতো কন্ট্রোল বসিয়ে বিস্তারিত নকশা করা যায়। সরল এন্ট্রির জন্য Wizard সুবিধাজনক; বিশেষ বিন্যাসের জন্য Design View প্রয়োজন হতে পারে।

Form Wizard খুলে Table/Query নির্বাচন করি, প্রয়োজনীয় ফিল্ড যুক্ত করি, layout বেছে নাম দিয়ে Finish করি। যেমন Student থেকে Name, BirthDate, GroupName ও Mobile নিয়ে StudentEntry তৈরি করা যায়। সব ফিল্ড ব্যবহারকারীকে দেখানো জরুরি নয়; পরিচিতি নম্বর স্বয়ংক্রিয় হলে তা শুধু প্রদর্শন করা যেতে পারে।

\textbf{Form View}-তে ব্যবহারকারী তথ্য দেখে ও অনুমতিসাপেক্ষে এন্ট্রি করে। \textbf{Layout View}-তে বাস্তব ডেটা দেখতে দেখতে অনেক বিন্যাস পরিবর্তন করা যায়। \textbf{Design View}-তে কন্ট্রোল, section এবং property বিস্তারিতভাবে বদলানো যায়। একটি কন্ট্রোলের প্রস্থ বাড়ানো নকশার পরিবর্তন; সেই ঘরে লেখা মোবাইল নম্বর বদলানো ডেটার পরিবর্তন।

\subsection{ফর্মের সেকশন}
\textbf{Form Header}-এ সাধারণত শিরোনাম বা নির্দেশনা থাকে। \textbf{Detail Section}-এ প্রতিটি রেকর্ডের ফিল্ড দেখানো ও এন্ট্রির কন্ট্রোল থাকে। \textbf{Form Footer}-এ সারাংশ, সহায়ক বোতাম বা অন্যান্য তথ্য রাখা যায়। যেমন StudentEntry-এর Header-এ কলেজের নাম, Detail-এ নাম ও জন্মতারিখ এবং Footer-এ Close বোতাম রাখা যায়।

ফর্ম দিয়ে নতুন রেকর্ড লিখলে বৈধ মান টেবিলে সংরক্ষণ হয়। বিদ্যমান রেকর্ডে মান বদলালে অনুমতি ও নিয়ম অনুযায়ী সংশোধন হয়। Validation Rule, Required এবং তালিকাভিত্তিক নির্বাচন এন্ট্রির ভুল কমায়। তবে উপযুক্ত রেঞ্জের ভেতরেও ভুল ব্যক্তির নম্বর লেখা সম্ভব; সুন্দর ফর্ম নিজে নির্ভুলতার নিশ্চয়তা নয়।

\subsection{মাস্টার ফর্ম ও সাবফর্ম}
এক-অনেক সম্পর্কের তথ্য দেখাতে \textbf{Master Form ও Subform} উপযোগী। মূল ফর্মে একজন শিক্ষার্থীর পরিচয় এবং সাবফর্মে তার একাধিক পরীক্ষার ফল দেখানো যায়। মূল ও উপফর্ম StudentID দিয়ে যুক্ত হলে বর্তমান শিক্ষার্থীর সংশ্লিষ্ট ফলগুলো দেখা যায়। সংযোগ ভুল হলে অন্য শিক্ষার্থীর ফলও দেখাতে পারে।

একাধিক টেবিলের ফিল্ড নিয়ে Wizard ব্যবহার করার আগে সম্পর্ক ও কী ঠিক থাকা প্রয়োজন। Form Design-এ Subform Control বসিয়ে Link Master Fields ও Link Child Fields উপযুক্তভাবে নির্ধারণ করা যায়। ফর্মে নাম দেখালেও সংযোগের ভিত্তি সাধারণত স্থিতিশীল পরিচিতি নম্বর; একই নামের দুই শিক্ষার্থীকে শুধু নাম দিয়ে যুক্ত করা অনির্ভরযোগ্য।

\subsection{ফর্মের বিন্যাস ও সম্পাদনার উদাহরণ}
একটি Columnar-ধরনের বিন্যাসে একজনের ফিল্ডগুলো একটির নিচে আরেকটি দেখা সুবিধাজনক; ভর্তি বা সদস্য এন্ট্রিতে এটি কার্যকর। Tabular বা Datasheet-ধরনের বিন্যাসে বহু রেকর্ডের তথ্য সারিভিত্তিক দেখা যায়। নির্বাচনের নাম ও উপলব্ধ layout Access-এর সংস্করণভেদে বদলাতে পারে, কিন্তু এক রেকর্ড বনাম বহু রেকর্ড প্রদর্শনের প্রয়োজন একই থাকে।

Form Wizard-এ উৎস বাছাইয়ের পর Selected Fields-এ প্রয়োজনীয় ফিল্ড নিতে হয়। সব ফিল্ড নির্বাচন না করলে বাদ দেওয়া ফিল্ড টেবিল থেকে মুছে যায় না। ফর্মের নাম StudentEntry এবং উৎস টেবিল Student হতে পারে। তৈরি হওয়া ফর্ম পরে Design View-তে খুলে Label, Text Box-এর মাপ ও অবস্থান বদলানো যায়।

\textbf{রেকর্ড সম্পাদনা:} নামের বানান সংশোধনের আগে সঠিক StudentID-এর রেকর্ডে আছ কি না দেখো। পরের রেকর্ডে গিয়ে ফিরে এসে মান যাচাই করো। নতুন শিক্ষার্থী এন্ট্রির জন্য নতুন রেকর্ডে যেতে হবে; বিদ্যমান শিক্ষার্থীর ঘরগুলো মুছে নতুন নাম দিলে আগের রেকর্ডই বদলে যেতে পারে। তাই navigation ও পরিচিতি যাচাই জরুরি।

\subsection{কন্ট্রোল নির্বাচন করে সমস্যার সমাধান}
একটি ফর্মে নাম, বিভাগ, ফি পরিশোধের অবস্থা ও রিপোর্ট খোলার সুবিধা দরকার। Name-এ Text Box, বিভাগে Combo Box, Paid-এ Check Box এবং রিপোর্ট খোলায় Command Button নির্বাচন করা যায়। প্রতিটি কন্ট্রোলের পাশে Label ব্যবহার করলে কী তথ্য দিতে হবে বোঝা যায়। শুধু সুন্দর রং বা ঘরের সংখ্যা দেখে কন্ট্রোল নির্বাচন নয়, তথ্যের ধরন ও ব্যবহার দেখতে হয়।
\begin{flushleft}\small
\begin{tabularx}{\linewidth}{|p{34mm}|X|}\hline
\chSixHeader পরিস্থিতি & প্রত্যাশিত আচরণ\\\hline
একই বিভাগ তালিকা থেকে নির্বাচন & Combo Box-এর বৈধ বিকল্প সংরক্ষণ\\\hline
Paid চিহ্নিত করা & সংশ্লিষ্ট Yes/No মান প্রদর্শন বা বদল\\\hline
শুধু নির্দেশনা দেখানো & Label; নতুন ডেটা-রেকর্ড নয়\\\hline
শিক্ষার্থী বদলানো & সাবফর্মে সংশ্লিষ্ট StudentID-এর ফল\\\hline
\end{tabularx}
\end{flushleft}

\chSixInterface{ch06_access_form_view.png}{Form View-এর পরিচিতি}{Label ফিল্ডের পরিচয় দেখায়; Text Box ও অন্যান্য Control দিয়ে রেকর্ডের মান দেখা বা পরিবর্তন করা যায়।}
